% Build twice for the table of contents:
%   pdflatex Digital_Twin_Divergence_Textbook.tex
%   pdflatex Digital_Twin_Divergence_Textbook.tex
\documentclass[11pt,oneside]{book}

\usepackage[margin=1in]{geometry}
\usepackage{amsmath,amssymb,bm}
\usepackage{booktabs,longtable,tabularx}
\usepackage{enumitem}
\usepackage{xcolor}
\usepackage{hyperref}
\usepackage{listings}
\usepackage{fancyhdr}

\hypersetup{
  colorlinks=true,
  linkcolor=blue!55!black,
  citecolor=blue!55!black,
  urlcolor=blue!55!black,
  pdftitle={Digital-Twin Localization, Validation, and Security},
  pdfauthor={Shreyas Udaya}
}

\pagestyle{fancy}
\fancyhf{}
\fancyhead[L]{Digital-Twin Localization, Validation, and Security}
\fancyhead[R]{\thepage}
\setlength{\headheight}{14pt}
\setlist[itemize]{leftmargin=*,itemsep=2pt}
\setlist[enumerate]{leftmargin=*,itemsep=2pt}
\setlength{\parindent}{0pt}
\setlength{\parskip}{6pt}

\definecolor{codegray}{RGB}{245,247,248}
\lstset{
  basicstyle=\ttfamily\small,
  backgroundcolor=\color{codegray},
  frame=single,
  breaklines=true,
  columns=fullflexible
}

\newcommand{\vect}[1]{\bm{#1}}
\newcommand{\E}{\mathbb{E}}
\newcommand{\R}{\mathbb{R}}
\newcommand{\projectpath}[1]{\texttt{#1}}
\newcommand{\keyidea}[1]{%
  \begin{center}
  \fbox{\parbox{0.92\textwidth}{\textbf{Key idea.} #1}}
  \end{center}}
\newcommand{\warningbox}[1]{%
  \begin{center}
  \fbox{\parbox{0.92\textwidth}{\textbf{Methodological warning.} #1}}
  \end{center}}

\title{\textbf{Digital-Twin Localization, Validation, and Security}\\
\large A Project-Specific Textbook for the UGV01, BN220 GNSS,\\
Tracked-Drive Odometry, Edge Telemetry, and AprilTag Validation}
\author{Shreyas Udaya}
\date{July 2026}

\begin{document}
\frontmatter
\maketitle

\chapter*{Purpose of This Book}
\addcontentsline{toc}{chapter}{Purpose of This Book}

This book is a guided foundation for one research project: a security-aware
digital twin of a Waveshare UGV01 tracked rover. The rover reports encoder,
inertial, GNSS, power, and timing telemetry to an edge computer. The edge
software predicts rover motion, fuses GNSS into an operational estimate, and
tests whether GNSS remains consistent with independent motion and transport
evidence.

The book has three goals:
\begin{enumerate}
  \item explain the mathematics well enough to derive and criticize the model;
  \item separate what the existing data establishes from what still requires
  independent physical ground truth; and
  \item connect every major idea to the repository implementation.
\end{enumerate}

This is an original synthesis of the sources in the bibliography. It is not a
replacement for the cited books and papers. Read the indicated primary sources
when a chapter becomes central to a paper claim.

\chapter*{How to Study}
\addcontentsline{toc}{chapter}{How to Study}

Read Chapters 1--7 before changing the estimator. Read Chapters 8--10 before
collecting ground-truth data. Read Chapters 11--12 before resuming the attack
campaign. For each chapter, answer:
\begin{enumerate}
  \item What quantity is estimated?
  \item Which observations constrain it?
  \item Which assumptions make the equations valid?
  \item Which metric tests the claim?
  \item What observation would falsify the claim?
\end{enumerate}

\tableofcontents

\mainmatter

\chapter{The Research Problem}

\section{The physical and digital systems}

The physical system consists of a tracked UGV01 rover, drive encoders, an IMU,
a BN220 GNSS receiver, an ESP32 firmware path, Wi-Fi transport, and an edge
computer. The digital system contains:
\begin{itemize}
  \item a tracked-drive motion model;
  \item an operational extended Kalman filter (EKF);
  \item a GNSS-independent security predictor;
  \item uncertainty and evidence-gating policies;
  \item innovation-based alarms;
  \item packet timing, loss, staleness, and queue instrumentation; and
  \item replay and statistical-analysis tools.
\end{itemize}

NIST describes credible digital twins in terms of verification, validation,
and uncertainty quantification (VVUQ), not visual resemblance alone
\cite{shao2023credibility,nistDT}. For this project:
\begin{itemize}
  \item \textbf{verification} asks whether the equations and software are
  implemented as intended;
  \item \textbf{validation} asks whether the model represents the rover
  accurately enough for its stated purpose; and
  \item \textbf{uncertainty quantification} asks whether the reported
  confidence reflects observed error.
\end{itemize}

\section{Why GNSS is used}

Encoders and an IMU provide relative motion. Encoder scale error and track slip
accumulate in position; gyro bias accumulates in heading. GNSS provides an
absolute position anchor outdoors, but can be degraded by obstruction,
reflection, multipath, interference, or adversarial manipulation
\cite{groves2013,psiaki2016}. GNSS is therefore both:
\begin{enumerate}
  \item an operational sensor that can bound dead-reckoning drift; and
  \item the primary untrusted measurement and attack surface.
\end{enumerate}

AprilTag camera tracking has a different role. It is an experimental reference
used to evaluate physical fidelity. It is not required for normal rover
deployment.

\keyidea{GNSS is an estimator input and security subject. AprilTag is the
independent experimental reference. These roles must never be interchanged.}

\section{Three distinct questions}

\begin{longtable}{p{0.23\textwidth}p{0.31\textwidth}p{0.36\textwidth}}
\toprule
Question & Suitable evidence & Unsuitable substitute\\
\midrule
\endhead
Does the software execute correctly? & Unit tests, replay equivalence, schema
checks, deterministic outputs & A visually plausible route\\
Is uncertainty calibrated? & NIS, NEES with ground truth, innovation bias and
whiteness & Small posterior-to-GNSS distance\\
Is the trajectory physically accurate? & Time-synchronized independent
position and heading ground truth & GNSS used by the same EKF\\
Does the detector resist attacks? & Frozen benign policy, matched attacks,
detection and false-alarm statistics & One successful attack example\\
\bottomrule
\end{longtable}

\chapter{Frames, State, and Tracked-Drive Motion}

\section{Coordinate frames}

Use explicit frames:
\begin{itemize}
  \item $\mathcal{G}$: geodetic latitude and longitude;
  \item $\mathcal{L}$: local planar east--north frame;
  \item $\mathcal{B}$: rover body frame;
  \item $\mathcal{C}$: camera frame; and
  \item $\mathcal{W}$: AprilTag ground-truth world frame.
\end{itemize}

At a local origin $(\phi_0,\lambda_0)$, a short-range approximation is
\begin{align}
x_G &\approx R_E\cos(\phi_0)(\lambda-\lambda_0),\\
y_G &\approx R_E(\phi-\phi_0),
\end{align}
where angles are in radians. This is adequate over a small test area, provided
the conversion is unit-tested and the origin is recorded.

\section{State and controls}

The planar state is
\[
\vect{x}_k =
\begin{bmatrix}
x_k & y_k & \theta_k
\end{bmatrix}^{T},
\]
where $(x,y)$ is local position and $\theta$ is heading. Let left and right
track travel over one interval be $\Delta s_L$ and $\Delta s_R$. With effective
track width $b$,
\begin{align}
\Delta s &= \frac{\Delta s_R+\Delta s_L}{2},\\
\Delta\theta_{\mathrm{enc}} &= \frac{\Delta s_R-\Delta s_L}{b},\\
v_k &= \frac{\Delta s}{\Delta t},\qquad
\omega_{\mathrm{enc},k}=\frac{\Delta\theta_{\mathrm{enc}}}{\Delta t}.
\end{align}

Encoder counts become distance through
\[
\Delta s_L=c_L\Delta n_L,\qquad
\Delta s_R=c_R\Delta n_R,
\]
where $c_L$ and $c_R$ are metres per count. These are \emph{effective}
parameters: track deformation and contact mechanics mean they need not equal a
simple wheel-circumference calculation.

\section{Nonlinear prediction}

A midpoint integration model is
\begin{align}
x_{k+1} &= x_k + v_k\Delta t
\cos\left(\theta_k+\frac{\omega_k\Delta t}{2}\right),\\
y_{k+1} &= y_k + v_k\Delta t
\sin\left(\theta_k+\frac{\omega_k\Delta t}{2}\right),\\
\theta_{k+1} &= \operatorname{wrap}(\theta_k+\omega_k\Delta t).
\end{align}

The EKF requires the Jacobian
\[
\vect{F}_k=\left.\frac{\partial f}{\partial\vect{x}}\right|_{\hat{\vect{x}}_k}.
\]
For the midpoint model, the nontrivial derivatives are the position
derivatives with respect to heading. A software test should compare the
analytic Jacobian with finite differences.

\section{Tracked-drive limitations}

The no-slip model fails when:
\begin{itemize}
  \item one track slides during a pivot turn;
  \item left and right traction differ;
  \item the rover coasts after a stop command;
  \item the effective track width changes with surface and turn rate; or
  \item encoder counts represent motor motion that did not become ground
  displacement.
\end{itemize}

These effects motivate slip indicators, but an encoder--IMU disagreement is
not itself proof of slip. Either sensor may be wrong. Independent trajectory
ground truth is needed to identify physical motion error.

\chapter{Sensor Models and Failure Modes}

\section{Encoders}

Encoder observations are precise measurements of drivetrain rotation, not
direct measurements of ground displacement. Important errors include scale,
sign, quantization, missed counts, backlash, and slip. Repeated straight and
turn trials estimate systematic effects; varying surfaces reveal transfer
limitations.

\section{Gyroscope}

Model the measured yaw rate as
\[
\omega_{\mathrm{imu},k}=\omega_k+b_g+n_{g,k},
\]
where $b_g$ is bias and $n_g$ is noise. A stationary prefix estimates
\[
\hat b_g=\operatorname{median}
\{\omega_{\mathrm{imu},k}:k\in\mathcal{I}_{\mathrm{stationary}}\}.
\]
The corrected rate is $\tilde\omega_{\mathrm{imu},k}
=\omega_{\mathrm{imu},k}-\hat b_g$.

The current implementation conservatively fuses rates as
\[
\omega_k=(1-\alpha)\omega_{\mathrm{enc},k}
+\alpha\tilde\omega_{\mathrm{imu},k},
\]
with a small frozen $\alpha$. A small weight is appropriate when encoder and
IMU turn rates have only moderate agreement.

\section{GNSS}

A position observation is
\[
\vect{z}^{G}_k=\vect{H}\vect{x}_k+\vect{v}_k,\qquad
\vect{H}=
\begin{bmatrix}
1&0&0\\0&1&0
\end{bmatrix}.
\]
A common model assumes $\vect{v}_k\sim\mathcal{N}(\vect{0},\vect{R}_k)$.
Real GNSS error can be biased, correlated, non-Gaussian, and environment
dependent. HDOP describes satellite geometry; it does not directly guarantee
position accuracy, especially under multipath.

For the indoor repeated-square logs, raw GNSS points occupy spatial extents far
larger than the commanded route. That observation establishes poor GNSS route
conformance, but not the true error of each fix because physical ground truth
was not recorded.

\section{Timing as a sensor}

The telemetry schema records source sample, send, edge arrival, queue,
estimate, and alarm times. Timing evidence can identify stale packets,
sequence gaps, delay, jitter, and queue buildup. Under the primary threat
model, these edge-observed fields are protected from a coordinate-only GNSS
attacker and can therefore help govern adaptation.

\chapter{Kalman Filtering and the EKF}

\section{Bayesian recursion}

Filtering estimates a posterior distribution
\[
p(\vect{x}_k\mid\vect{z}_{1:k},\vect{u}_{1:k}).
\]
The Kalman filter is the exact Gaussian solution for linear dynamics and
measurements \cite{kalman1960,welchbishop}. The EKF applies the same recursion
after local linearization \cite{thrun2005,jazwinski1970}.

\section{Prediction}

Given posterior $(\hat{\vect{x}}_{k-1}^{+},\vect{P}_{k-1}^{+})$,
\begin{align}
\hat{\vect{x}}_k^{-} &= f(\hat{\vect{x}}_{k-1}^{+},\vect{u}_k),\\
\vect{P}_k^{-} &= \vect{F}_k\vect{P}_{k-1}^{+}\vect{F}_k^T+\vect{Q}_k.
\end{align}

$\vect{Q}_k$ represents uncertainty in the process model over the interval. It
is not an arbitrary smoothing knob. Increasing it says that motion prediction
is less trustworthy.

\section{Measurement update}

For GNSS measurement $\vect{z}_k$,
\begin{align}
\vect{\nu}_k &= \vect{z}_k-\vect{H}\hat{\vect{x}}_k^{-},\\
\vect{S}_k &= \vect{H}\vect{P}_k^{-}\vect{H}^T+\vect{R}_k,\\
\vect{K}_k &= \vect{P}_k^{-}\vect{H}^T\vect{S}_k^{-1},\\
\hat{\vect{x}}_k^{+} &= \hat{\vect{x}}_k^{-}+\vect{K}_k\vect{\nu}_k.
\end{align}

The innovation $\vect{\nu}_k$ is a \emph{pre-update} discrepancy. A
post-update EKF--GNSS residual is partly circular because the same GNSS sample
has already moved the EKF toward itself.

\section{Meaning of the covariance matrices}

\begin{itemize}
  \item $\vect{P}$: uncertainty in the estimated state.
  \item $\vect{Q}$: uncertainty introduced by motion prediction.
  \item $\vect{R}$: uncertainty in the measurement.
  \item $\vect{S}$: expected covariance of the innovation.
\end{itemize}

Scaling $\vect{P}$, $\vect{Q}$, and $\vect{R}$ by the same positive scalar
leaves the Kalman gain unchanged. It can change normalized consistency scores
without changing the state trajectory. Therefore, covariance calibration and
trajectory improvement are different tasks.

\warningbox{Never claim that improved NIS coverage proves improved position
accuracy. The covariance may improve while the trajectory remains identical.}

\chapter{Consistency: Innovation, NIS, and NEES}

\section{Normalized innovation squared}

The normalized innovation squared is
\[
d_k=\vect{\nu}_k^T\vect{S}_k^{-1}\vect{\nu}_k.
\]
Under a correct linear-Gaussian model with two position measurements,
$d_k\sim\chi_2^2$. The 95th percentile is approximately $5.991$. Thus, over
many independent representative updates, approximately 95\% should fall below
that threshold \cite{brumback1987}.

Interpretation:
\begin{itemize}
  \item too many large NIS values suggest underestimated uncertainty,
  unmodelled bias, outliers, or model error;
  \item unusually small NIS values can suggest overestimated covariance; and
  \item correct coverage alone does not establish whiteness, independence, or
  physical accuracy.
\end{itemize}

\section{Innovation mean and whiteness}

A credible model should have approximately zero-mean innovations:
\[
\E[\vect{\nu}_k]\approx\vect{0}.
\]
It should also avoid strong temporal correlation. Persistent innovation bias
means the model repeatedly predicts on one side of the observation. Serial
correlation means information remains in the residual that the model has not
explained.

\section{NEES}

With independent ground truth $\vect{x}^{GT}_k$, define state error
\[
\vect{e}_k=\hat{\vect{x}}_k-\vect{x}^{GT}_k.
\]
The normalized estimation error squared is
\[
\mathrm{NEES}_k=\vect{e}_k^T\vect{P}_k^{-1}\vect{e}_k.
\]
NEES evaluates state-space consistency and requires ground truth. NIS operates
in measurement space and does not.

\section{Current project interpretation}

The current benign corpus yields approximately 91.1\% run-averaged NIS
coverage after development-only covariance scaling, compared with a nominal
95\% target. Development coverage is closer to nominal than held-out
validation and test coverage. This means uncertainty calibration is promising
but does not fully transfer across conditions. It is not ``91.1\% localization
accuracy.''

\chapter{A Two-Branch Digital Twin}

\section{Operational branch}

The operational EKF predicts from encoder/IMU motion and accepts GNSS
corrections. Its purpose is navigation-state estimation when GNSS is useful.
Because GNSS updates the state, the operational branch can be pulled by a
biased or attacked GNSS stream.

\section{Security-reference branch}

The security predictor begins from a clean alignment interval and then
propagates using encoder/IMU motion without GNSS state correction. GNSS is
scored against this prediction:
\[
\vect{\nu}^{S}_k=\vect{z}^{G}_k-\vect{H}\hat{\vect{x}}^{S,-}_k.
\]
This branch can drift naturally, but it avoids the direct circularity of using
the suspect coordinate to create the reference against which it is tested.

\section{Initialization}

Initial heading is weakly observable while stationary. The implementation uses
early motion to align a short dead-reckoned path with the initial GNSS
displacement pattern. This improves repeatability, but an initial GNSS bias can
still rotate or translate the local frame. AprilTag ground truth is needed to
evaluate the resulting physical heading.

\section{Evidence-gated uncertainty adaptation}

Adaptive covariance can be useful when motion uncertainty changes
\cite{mehra1970}. It becomes dangerous if an untrusted GNSS residual directly
increases the covariance that normalizes its own detector score. The causal
loop is
\[
\text{GNSS manipulation}\rightarrow\text{residual feature}
\rightarrow\vect{Q}\text{ inflation}\rightarrow\vect{S}\text{ inflation}
\rightarrow\text{smaller normalized score}.
\]

The evidence-gated design limits this path by:
\begin{itemize}
  \item excluding GNSS coordinate residuals from process features;
  \item using encoder, IMU, and edge timing as independent evidence;
  \item bounding covariance proposals; and
  \item freezing the previous accepted proposal when evidence is unhealthy.
\end{itemize}

\chapter{Physical Fidelity and Trajectory Metrics}

\section{Why GNSS agreement is insufficient}

Let $\hat{\vect{p}}_k$ be an EKF position and $\vect{z}^{G}_k$ the GNSS
position. The distance
\[
r_k=\|\hat{\vect{p}}_k-\vect{z}^{G}_k\|
\]
measures sensor/model agreement. It is not physical position error. If the EKF
consumes GNSS, both can share the same bias.

The physical error is
\[
e_k=\|\hat{\vect{p}}_k-\vect{p}^{GT}_k\|.
\]
It is unavailable until an independent reference supplies
$\vect{p}^{GT}_k$.

\section{Absolute trajectory error}

After one fixed, documented frame registration, absolute trajectory error
(ATE) is computed over synchronized poses \cite{sturm2012}:
\[
\mathrm{ATE}_{\mathrm{RMSE}}=
\sqrt{\frac{1}{N}\sum_{k=1}^{N}
\|\hat{\vect{p}}_k-\vect{p}^{GT}_k\|^2}.
\]
Report RMSE, median, 95th percentile, maximum, and percentages within
pre-registered tolerances. Do not optimize a separate alignment for every run
if global position accuracy is part of the claim.

\section{Relative pose error}

Relative pose error (RPE) compares estimated and true changes over an interval
$\Delta$. It measures local drift even when global alignment differs:
\[
\vect{E}_{k,\Delta}=
\left((\vect{T}^{GT}_k)^{-1}\vect{T}^{GT}_{k+\Delta}\right)^{-1}
\left((\hat{\vect{T}}_k)^{-1}\hat{\vect{T}}_{k+\Delta}\right).
\]
Report translational and rotational components at meaningful time or distance
intervals.

\section{Route-shape metrics}

For controlled routes, additionally report:
\begin{itemize}
  \item cross-track RMSE and p95;
  \item heading mean absolute error and p95;
  \item straight-segment curvature;
  \item side-length and corner-angle errors;
  \item loop-closure error;
  \item path-length ratio; and
  \item Fr\'echet or time-aligned trajectory distance.
\end{itemize}

A commanded 0.5-m square is a protocol, not ground truth. Slip and imperfect
control mean the rover may not execute the command exactly.

\section{What ``good enough'' means}

Accuracy is application-relative. Define acceptance before final collection:
\begin{itemize}
  \item reference-system uncertainty must be substantially below the smallest
  estimator error of interest;
  \item estimator error must be small relative to mission tolerance and attack
  magnitudes;
  \item heading and local motion must be accurate enough that benign
  dead-reckoning does not resemble an attack; and
  \item performance must transfer across held-out surfaces and speeds.
\end{itemize}

\chapter{AprilTag Ground Truth}

\section{Why fiducials}

AprilTag provides identifiable planar markers whose corners support robust
pose estimation \cite{olson2011apriltag,wang2016apriltag2}. A calibrated fixed
camera can track a rigid rover-mounted tag while surveyed stationary tags
define the world frame.

\section{Camera calibration}

The pinhole model is
\[
s
\begin{bmatrix}u\\v\\1\end{bmatrix}
=
\vect{K}
\begin{bmatrix}\vect{R}&\vect{t}\end{bmatrix}
\begin{bmatrix}X\\Y\\Z\\1\end{bmatrix},
\]
where $\vect{K}$ contains focal lengths and principal point. Lens distortion
must be estimated. Zhang's planar-pattern method obtains intrinsics from
multiple pattern orientations and refines them nonlinearly
\cite{zhang2000calibration}.

A ChArUco board helps calibrate the camera. It is not the trajectory reference
by itself. Record the printed square dimensions actually measured, not merely
the design dimensions.

\section{Planar world mapping}

For a rover moving on a planar floor, a homography maps undistorted image
coordinates to world-plane coordinates:
\[
s\vect{p}_{\mathrm{image}}=\vect{H}_{CW}\vect{p}_{\mathrm{world}}.
\]
Use at least four non-collinear surveyed world references distributed around
the route. Validate the mapping on measured points not used to fit it.

\section{Time synchronization}

Video and telemetry have different clocks and sample rates. The processing
pipeline must:
\begin{enumerate}
  \item preserve original frame timestamps;
  \item estimate camera-to-telemetry offset using a visible and
  telemetry-observable synchronization event;
  \item interpolate one trajectory onto the other time base;
  \item quantify residual synchronization error; and
  \item exclude long tag-occlusion gaps rather than hiding them.
\end{enumerate}

\section{Reference uncertainty}

AprilTag is not automatically perfect ground truth. Quantify:
\begin{itemize}
  \item calibration reprojection error;
  \item static position and heading repeatability;
  \item error at held-out measured floor points;
  \item sensitivity to viewing angle and distance;
  \item detection coverage; and
  \item timestamp uncertainty.
\end{itemize}

\keyidea{A reference becomes credible through calibration and validation, not
because it uses a particular marker technology.}

\chapter{Robust GNSS Use}

\section{Quality-aware measurement covariance}

When GNSS quality degrades, increasing $\vect{R}_k$ reduces the Kalman gain and
therefore the measurement's influence. Quality features may include satellite
count, HDOP, innovation history, environment classification, and receiver
status. HDOP alone cannot identify coherent multipath bias.

\section{Innovation gating}

A hard gate rejects updates when NIS exceeds a threshold. This protects
against large outliers, but can reject legitimate measurements after the
predictor has drifted. A gate threshold must be selected on benign development
data and evaluated on untouched benign runs.

\section{Huber-style robustification}

Huber methods reduce the weight of large residuals rather than fully rejecting
them \cite{gandhi2010}. They can improve outlier resistance but do not solve
slow coherent bias automatically.

\section{Sequential detection}

CUSUM accumulates persistent evidence:
\[
C_k=\max(0,C_{k-1}+s_k-\kappa),
\]
where $s_k$ is a residual score and $\kappa$ is an allowance. An alarm occurs
when $C_k>h$. It trades detection delay against false alarms and is useful for
slow changes \cite{page1954,murguia2016}.

\section{A correct policy hierarchy}

For this project, a defensible hierarchy is:
\begin{enumerate}
  \item fuse GNSS normally when independent evidence and innovations are
  healthy;
  \item down-weight uncertain but plausible GNSS;
  \item reject clearly inconsistent updates;
  \item continue the independent security prediction;
  \item alarm only under a frozen persistence policy; and
  \item expose degraded mode rather than silently claiming accuracy.
\end{enumerate}

\chapter{Threat Models and Attack Evaluation}

\section{Threat boundary}

The primary attacker modifies GNSS coordinate content delivered to the edge
pipeline. The attacker does not compromise encoders, IMU, edge timestamps,
alarm code, commands, or the AprilTag reference. This makes independent
evidence meaningful.

Semantic coordinate injection evaluates estimator-visible consequences but is
not an over-the-air RF spoofing experiment. GNSS spoofing literature motivates
the threat, but claims must remain scoped to the implemented attack interface
\cite{psiaki2016,shen2020}.

\section{Attack families}

\begin{itemize}
  \item \textbf{Step bias}: sudden coordinate displacement.
  \item \textbf{Slow drift}: gradually increasing displacement.
  \item \textbf{Freeze}: repeated old coordinate.
  \item \textbf{Replay}: previously observed coordinate sequence.
  \item \textbf{Delivery perturbation}: delay, jitter, loss, or reordering,
  evaluated separately from coordinate manipulation.
\end{itemize}

\section{Baselines}

Compare against mechanisms with distinct assumptions:
\begin{itemize}
  \item raw GNSS jump and raw model residual;
  \item fixed-covariance NIS;
  \item innovation-gated EKF;
  \item Huber EKF;
  \item CUSUM on whitened innovations;
  \item residual-coupled adaptive covariance;
  \item GNSS-independent adaptive covariance; and
  \item evidence-gated adaptation.
\end{itemize}

Secure-state-estimation research emphasizes explicit attack sparsity,
observability, and trust assumptions \cite{pasqualetti2013,fawzi2014}. A method
is not secure merely because it uses multiple sensors.

\section{Evaluation quantities}

Report:
\begin{itemize}
  \item benign run-level false-alarm probability;
  \item detection probability with run-clustered confidence intervals;
  \item time and distance to first alarm;
  \item maximum ground-truth deviation before alarm;
  \item tolerance-exceeding divergence before alarm;
  \item attack magnitude for 50\%, 90\%, and 95\% detection; and
  \item compute and transport latency.
\end{itemize}

Without physical ground truth, call a metric ``paired reported-state
divergence,'' not physical harm.

\chapter{Experimental Design}

\section{Separate the corpora}

\begin{enumerate}
  \item \textbf{Development}: fit geometry, covariance, and thresholds.
  \item \textbf{Validation}: choose among a small number of predefined
  alternatives.
  \item \textbf{Prospective test}: untouched evaluation after freezing.
\end{enumerate}

Runs, not individual telemetry rows, are the primary independent units.
Bootstrap complete runs to preserve within-run temporal dependence.

\section{Benign protocol}

Use:
\begin{itemize}
  \item stationary prefixes for gyro and GNSS characterization;
  \item measured straight runs for scale;
  \item turns for effective track width and heading;
  \item squares for corners and closure;
  \item figure-eights for both turn directions; and
  \item an irregular fixed route for transfer beyond a square.
\end{itemize}

Balance speed, surface, direction, and GNSS environment. Indoor data should be
labeled degraded GNSS and analyzed separately from outdoor nominal GNSS.

\section{Run validity}

Predefine valid-run requirements: complete telemetry schema, successful time
alignment, sufficient tag coverage, no camera movement, acceptable reference
quality, and complete route execution. Preserve invalid runs and reasons.

\section{Avoiding leakage}

Do not:
\begin{itemize}
  \item tune on the prospective set;
  \item exclude a difficult benign run because it raises false alarms;
  \item select attack injection times after viewing detector scores;
  \item fit a different frame alignment to make each result look better; or
  \item count thousands of detector evaluations as thousands of independent
  physical experiments.
\end{itemize}

\chapter{Reading the Repository}

\section{Core modules}

\begin{longtable}{p{0.40\textwidth}p{0.50\textwidth}}
\toprule
Path & Question to answer\\
\midrule
\endhead
\projectpath{DigitalTwin/kinematics.py} & How do encoder counts become
$v$ and $\omega$?\\
\projectpath{DigitalTwin/motion.py} & How are gyro bias, yaw fusion,
initial heading, and slip indicators calculated?\\
\projectpath{DigitalTwin/ekf.py} & Where are prediction, innovation,
gain, and update implemented?\\
\projectpath{DigitalTwin/uncertainty.py} & Which features influence
$\vect{Q}$ and $\vect{R}$?\\
\projectpath{DigitalTwin/security.py} & What remains independent of GNSS,
and what opens or closes the evidence gate?\\
\projectpath{DigitalTwin/detector.py} & Which score and threshold generate
an instantaneous detection?\\
\projectpath{DigitalTwin/alarm.py} & How does persistence convert detections
into a run-level alarm?\\
\projectpath{DigitalTwin/analysis/real\_data\_study.py} & How are physical
runs, attacks, detector variants, and outputs organized?\\
\projectpath{DigitalTwin/analysis/digital\_twin\_accuracy.py} & Which reported
quantities are sensor agreement, consistency, or physical accuracy?\\
\bottomrule
\end{longtable}

\section{Code-reading exercise}

For one benign log:
\begin{enumerate}
  \item identify the stationary interval and gyro-bias estimate;
  \item calculate one encoder-derived $v$ and $\omega$ by hand;
  \item trace one EKF prediction;
  \item calculate the GNSS innovation and $\vect{S}$;
  \item calculate NIS and compare with 5.991;
  \item determine whether the evidence gate accepts adaptation; and
  \item explain why none of those calculations establishes physical accuracy.
\end{enumerate}

\chapter{Interpreting the Current Evidence}

\section{What is established}

The project currently demonstrates:
\begin{itemize}
  \item real UGV01 telemetry acquisition and replay;
  \item GNSS, encoder, IMU, and timing instrumentation;
  \item deterministic attack injection and comparator evaluation;
  \item an operational EKF and GNSS-independent security predictor;
  \item gyro-bias correction and conservative encoder/IMU yaw fusion;
  \item development/validation/test separation for recent covariance work; and
  \item internal consistency and sensor-agreement reports.
\end{itemize}

\section{What current numbers mean}

The current operational EKF--GNSS pooled RMSE is approximately 1.297 m across
the existing evaluated samples. Outdoor rough-surface agreement is better than
indoor smooth-floor agreement. These are discrepancies with GNSS, not physical
position errors. The indoor GNSS trace extent is much larger than the
commanded square, so indoor runs should be treated as degraded-GNSS stress
tests.

Current NIS coverage of roughly 91.1\% indicates imperfect but informative
uncertainty calibration. It does not quantify localization accuracy.

\section{What remains unknown}

Until synchronized AprilTag data exist, the project cannot state:
\begin{itemize}
  \item physical ATE or RPE;
  \item physical heading or cross-track error;
  \item whether the operational or security trajectory is closer to reality;
  \item NEES calibration; or
  \item physical deviation before an attack alarm.
\end{itemize}

\section{Correct next milestone}

The next scientific milestone is not a larger replay campaign. It is:
\begin{enumerate}
  \item calibrate and validate the camera reference;
  \item collect pilot trajectories with synchronized AprilTag ground truth;
  \item correct and freeze the benign estimator;
  \item collect an untouched prospective benign corpus;
  \item demonstrate fidelity and false-alarm behavior; and
  \item only then run the confirmatory attack evaluation.
\end{enumerate}

\chapter{Review Questions}

\begin{enumerate}
  \item Why can an EKF have good GNSS agreement but poor physical accuracy?
  \item Why can NIS improve without any change in the state trajectory?
  \item What makes an evidence source independent under the primary threat
  model?
  \item Why is a commanded square not ground truth?
  \item How do ATE and RPE answer different questions?
  \item What errors can an AprilTag reference introduce?
  \item Why must indoor and outdoor GNSS conditions be stratified?
  \item How can adaptive covariance reduce detector sensitivity?
  \item Why are telemetry rows not independent experimental trials?
  \item What evidence must exist before the phrase ``physical harm before
  alarm'' is defensible?
\end{enumerate}

\appendix

\chapter{Notation}

\begin{longtable}{p{0.20\textwidth}p{0.70\textwidth}}
\toprule
Symbol & Meaning\\
\midrule
\endhead
$\vect{x},\hat{\vect{x}}$ & true and estimated state\\
$\vect{P}$ & state-estimation covariance\\
$\vect{Q}$ & process-noise covariance\\
$\vect{R}$ & measurement-noise covariance\\
$\vect{F}$ & process-model Jacobian\\
$\vect{H}$ & measurement matrix or Jacobian\\
$\vect{\nu}$ & pre-update innovation\\
$\vect{S}$ & innovation covariance\\
$\vect{K}$ & Kalman gain\\
$d_k$ & normalized innovation squared\\
$v,\omega$ & forward and angular velocity\\
$b_g$ & gyro yaw-rate bias\\
$\vect{T}$ & homogeneous pose transformation\\
\bottomrule
\end{longtable}

\chapter{A One-Page Claim Checklist}

Before making a claim, complete:
\begin{enumerate}
  \item \textbf{Quantity}: What exactly is measured?
  \item \textbf{Reference}: Against what independent reference?
  \item \textbf{Units}: Metres, degrees, seconds, probability, or normalized
  score?
  \item \textbf{Population}: Updates, routes, or independent physical runs?
  \item \textbf{Split}: Development, validation, or prospective test?
  \item \textbf{Uncertainty}: What confidence interval and reference error?
  \item \textbf{Scope}: Indoor/outdoor, speed, surface, and threat assumptions?
  \item \textbf{Falsifier}: What result would make the claim fail?
\end{enumerate}

\backmatter

\begin{thebibliography}{99}

\bibitem{shao2023credibility}
G. Shao, J. Hightower, and W. Schindel,
``Credibility consideration for digital twins in manufacturing,''
\emph{Manufacturing Letters}, vol. 35, pp. 24--28, 2023.
\href{https://doi.org/10.1016/j.mfglet.2022.11.009}
{doi:10.1016/j.mfglet.2022.11.009}.

\bibitem{nistDT}
National Institute of Standards and Technology,
``Digital Twins for Advanced Manufacturing,'' accessed July 2026.
\url{https://www.nist.gov/programs-projects/digital-twins-advanced-manufacturing}.

\bibitem{kalman1960}
R. E. Kalman, ``A new approach to linear filtering and prediction problems,''
\emph{Journal of Basic Engineering}, vol. 82, no. 1, pp. 35--45, 1960.
\href{https://doi.org/10.1115/1.3662552}{doi:10.1115/1.3662552}.

\bibitem{welchbishop}
G. Welch and G. Bishop, ``An introduction to the Kalman filter,''
University of North Carolina at Chapel Hill, Tech. Rep. TR 95-041.
\url{https://www.cs.unc.edu/~welch/media/pdf/kalman_intro.pdf}.

\bibitem{thrun2005}
S. Thrun, W. Burgard, and D. Fox,
\emph{Probabilistic Robotics}. MIT Press, 2005.
\url{https://mitpress.mit.edu/9780262201629/probabilistic-robotics/}.

\bibitem{jazwinski1970}
A. H. Jazwinski, \emph{Stochastic Processes and Filtering Theory}.
Academic Press, 1970.

\bibitem{groves2013}
P. D. Groves,
\emph{Principles of GNSS, Inertial, and Multisensor Integrated Navigation
Systems}, 2nd ed. Artech House, 2013.

\bibitem{psiaki2016}
M. L. Psiaki and T. E. Humphreys, ``GNSS spoofing and detection,''
\emph{Proceedings of the IEEE}, vol. 104, no. 6, pp. 1258--1270, 2016.
\href{https://doi.org/10.1109/JPROC.2016.2526658}
{doi:10.1109/JPROC.2016.2526658}.

\bibitem{sturm2012}
J. Sturm, N. Engelhard, F. Endres, W. Burgard, and D. Cremers,
``A benchmark for the evaluation of RGB-D SLAM systems,''
in \emph{Proc. IEEE/RSJ IROS}, 2012, pp. 573--580.
\href{https://doi.org/10.1109/IROS.2012.6385773}
{doi:10.1109/IROS.2012.6385773}.

\bibitem{zhang2000calibration}
Z. Zhang, ``A flexible new technique for camera calibration,''
\emph{IEEE Transactions on Pattern Analysis and Machine Intelligence},
vol. 22, no. 11, pp. 1330--1334, 2000.
\href{https://doi.org/10.1109/34.888718}{doi:10.1109/34.888718}.

\bibitem{olson2011apriltag}
E. Olson, ``AprilTag: A robust and flexible visual fiducial system,''
in \emph{Proc. IEEE ICRA}, 2011, pp. 3400--3407.
\url{https://april.eecs.umich.edu/papers/details.php?name=olson2011tags}.

\bibitem{wang2016apriltag2}
J. Wang and E. Olson, ``AprilTag 2: Efficient and robust fiducial detection,''
in \emph{Proc. IEEE/RSJ IROS}, 2016, pp. 4193--4198.
\href{https://doi.org/10.1109/IROS.2016.7759617}
{doi:10.1109/IROS.2016.7759617}.

\bibitem{gandhi2010}
M. A. Gandhi and L. Mili,
``Robust Kalman filter based on a generalized maximum-likelihood-type
estimator,'' \emph{IEEE Transactions on Signal Processing}, vol. 58, no. 5,
pp. 2509--2520, 2010.
\href{https://doi.org/10.1109/TSP.2009.2039731}
{doi:10.1109/TSP.2009.2039731}.

\bibitem{brumback1987}
B. D. Brumback and M. D. Srinath,
``A chi-square test for fault-detection in Kalman filters,''
\emph{IEEE Transactions on Automatic Control}, vol. 32, no. 6,
pp. 552--554, 1987.
\href{https://doi.org/10.1109/TAC.1987.1104658}
{doi:10.1109/TAC.1987.1104658}.

\bibitem{mehra1970}
R. K. Mehra, ``On the identification of variances and adaptive Kalman
filtering,'' \emph{IEEE Transactions on Automatic Control}, vol. 15, no. 2,
pp. 175--184, 1970.
\href{https://doi.org/10.1109/TAC.1970.1099422}
{doi:10.1109/TAC.1970.1099422}.

\bibitem{page1954}
E. S. Page, ``Continuous inspection schemes,''
\emph{Biometrika}, vol. 41, no. 1/2, pp. 100--115, 1954.
\href{https://doi.org/10.2307/2333009}{doi:10.2307/2333009}.

\bibitem{murguia2016}
C. Murguia and J. Ruths,
``CUSUM and chi-squared attack detection of compromised sensors,''
in \emph{Proc. IEEE Conference on Control Applications}, 2016, pp. 474--480.
\href{https://doi.org/10.1109/CCA.2016.7587875}
{doi:10.1109/CCA.2016.7587875}.

\bibitem{pasqualetti2013}
F. Pasqualetti, F. D\"orfler, and F. Bullo,
``Attack detection and identification in cyber-physical systems,''
\emph{IEEE Transactions on Automatic Control}, vol. 58, no. 11,
pp. 2715--2729, 2013.
\href{https://doi.org/10.1109/TAC.2013.2266831}
{doi:10.1109/TAC.2013.2266831}.

\bibitem{fawzi2014}
H. Fawzi, P. Tabuada, and S. Diggavi,
``Secure estimation and control for cyber-physical systems under adversarial
attacks,'' \emph{IEEE Transactions on Automatic Control}, vol. 59, no. 6,
pp. 1454--1467, 2014.
\href{https://doi.org/10.1109/TAC.2014.2303233}
{doi:10.1109/TAC.2014.2303233}.

\bibitem{shen2020}
J. Shen, J. Y. Won, Z. Chen, and Q. A. Chen,
``Drift with Devil: Security of multi-sensor fusion based localization in
high-level autonomous driving under GPS spoofing,''
in \emph{Proc. 29th USENIX Security Symposium}, 2020, pp. 931--948.
\url{https://www.usenix.org/conference/usenixsecurity20/presentation/shen}.

\end{thebibliography}

\end{document}
